% !TEX root = ../main.tex
\section{Baseline comparison}
\label{sec:baselines}

The original work compares three methods at $N = 8$: standard-average FedNPG
\eqref{eq:stdfednpg}, FedNPG-ADMM, and first-order federated Proximal Policy Optimization (FedPPO) with clipping
parameter $0.2$. We implement all three in the same codebase, holding
everything except the server-side aggregation identical: same environments, normalization, GAE, critic and logged metrics. On Hopper all three
methods use the reduced configuration of Section~\ref{sec:implementation}.

\subsection{Standard-average FedNPG}

This arm solves \eqref{eq:stdfednpg} exactly by conjugate gradient on the
pooled Fisher, exploiting $\sum_i \bH_i v = N \cdot \bF_{\text{pooled}} v$.
The solve is matrix-free, but the \emph{metering} is not: the channel is
charged the full $d^2$ elements per agent per round.

The obvious objection is that a baseline meant to pay an $\Oh(d^2)$ cost ought
to pay it. Materializing $\sum_i \bH_i$ is impossible at Humanoid scale
(Table~\ref{tab:commcost}: \SI{2.1}{\peta\byte}), so a literal implementation
would not exist for the environment where the $\Oh(d^2)$ cost matters most.
The claim under test concerns \emph{uplink bytes}, not server FLOPs; the paper
measures ``the number of transmitted parameters with double precision in each
agent''. We therefore compute the identical direction cheaply and report the
true transmission cost. Without that metering the baseline would look as cheap as FedNPG-ADMM, and the
comparison would measure nothing.

\subsection{FedPPO}

The paper names FedPPO and gives its clipping parameter, but specifies neither
the aggregation scheme nor the local epoch count. We implement the
FedAvg reading: each agent runs local PPO epochs from the broadcast $\bthe$
and uplinks its parameter delta, and the server averages. The alternative, a
single clipped-surrogate gradient step per round, would parallel
Algorithm~\ref{alg:fednpgadmm}'s structure more neatly, but it strips PPO of
the multi-epoch minibatch reuse that its clipping exists to make safe, so a
comparison against it would not reflect PPO as it is normally run.

\subsection{Communication cost, measured}

\begin{table}[t]
  \centering
  \caption{Measured per-iteration uplink, Swimmer-v4 with four agents
  ($d = \num{4868}$), at 64 bits per parameter. Measured values match the closed
  form exactly.}
  \label{tab:measuredcomm}
  \begin{tabular}{llrr}
    \toprule
    Method & Transmitted & Bytes/iteration & Relative \\
    \midrule
    Standard-average FedNPG & $\bH_i$ ($d^2$) $+\ \bg_i$ & \num{758473344} & $2434\times$ \\
    FedNPG-ADMM             & $\by_i,\ \bg_i$ (2 vectors) & \num{311552}   & $2\times$ \\
    FedPPO                  & $\Delta\bthe_i$ (1 vector)  & \num{155776}   & $1\times$ \\
    \bottomrule
  \end{tabular}
\end{table}

These measurements let us check the paper's headline saving. Per agent per
round, FedNPG-ADMM transmits $\num{2434}$ times less than standard-average
FedNPG on Swimmer and $\num{363537}$ times less on Humanoid
(Table~\ref{tab:commcost}). \citet{lan2023fednpgadmm} report about six orders
of magnitude on Humanoid, which matches, and four on Swimmer, where we measure
$10^{3.4}$. Part of the gap is that FedNPG-ADMM sends two vectors per round,
so the saving is about $d/2$ rather than $d$; the paper does not say how its
own figure was computed.

Table~\ref{tab:measuredcomm} also shows that FedNPG-ADMM is not the cheapest
of the three methods, although the paper's framing invites a reader to assume
it is. It is $\Oh(d)$, but with twice the constant of
FedPPO, since it transmits two vectors per round where FedPPO transmits one.
The argument for the method therefore rests on what it achieves per round
rather than on what it costs per round, and a figure plotted against iterations
cannot show that. We consequently plot reward against \emph{cumulative bytes},
which is the axis on which the two effects can be read together.

\subsection{Results}

\begin{table}[t]
  \centering
  \caption{Final rewards at $N=8$, mean $\pm$ standard deviation across five
  seeds, computed as in Table~\ref{tab:reproduction}. Hopper uses the
  configuration of Section~\ref{sec:implementation} for all three methods.}
  \label{tab:baselines}
  \begin{tabular}{lrrr}
    \toprule
    Method & Swimmer-v4 & Hopper-v4 & Humanoid-v4 \\
    \midrule
    FedNPG-ADMM             & $128.4 \pm 0.8$ & $\mathbf{2656 \pm 286}$ & $\mathbf{6329 \pm 334}$ \\
    Standard-average FedNPG & $127.1 \pm 0.5$ & $2494 \pm 677$ & $4699 \pm 298$ \\
    FedPPO                  & $\mathbf{129.8 \pm 0.7}$ & $1571 \pm 588$ & $5780 \pm 158$ \\
    \bottomrule
  \end{tabular}
\end{table}

\begin{figure}[t]
  \centering
  \begin{subfigure}[b]{0.49\textwidth}
    %% Figure 3 analogue, Swimmer-v4: three methods at N=8, mean over five seeds,
%% band spans min to max seed.
%% Data: figures/baselines_swimmer.dat (scripts/export_figures.py).
\begin{tikzpicture}
\begin{axis}[
    width=0.96\linewidth,
    height=0.78\linewidth,
    xlabel={Iteration},
    xmin=1, xmax=1000,
    xtick distance=250,
    ylabel={Average batch return},
    ymin=0, ymax=140,
    ytick distance=35,
    ymajorgrids=true, xmajorgrids=true, grid style=dashed,
    tick label style={font=\footnotesize},
    label style={font=\footnotesize},
    legend style={font=\footnotesize, cells={anchor=west},
                  anchor=south east, at={(0.99,0.02)}},
]

\addplot[name path=admmilo, draw=none, forget plot]
    table[x=iteration, y=admm_lo] {figures/baselines_swimmer.dat};
\addplot[name path=admmihi, draw=none, forget plot]
    table[x=iteration, y=admm_hi] {figures/baselines_swimmer.dat};
\addplot[colorA, opacity=0.15, forget plot] fill between[of=admmilo and admmihi];
\addplot[colorA, line width=0.9pt]
    table[x=iteration, y=admm_mean] {figures/baselines_swimmer.dat};
\addlegendentry{FedNPG-ADMM};

\addplot[name path=stdilo, draw=none, forget plot]
    table[x=iteration, y=std_lo] {figures/baselines_swimmer.dat};
\addplot[name path=stdihi, draw=none, forget plot]
    table[x=iteration, y=std_hi] {figures/baselines_swimmer.dat};
\addplot[colorB, opacity=0.15, forget plot] fill between[of=stdilo and stdihi];
\addplot[colorB, line width=0.9pt]
    table[x=iteration, y=std_mean] {figures/baselines_swimmer.dat};
\addlegendentry{Standard FedNPG};

\addplot[name path=ppoilo, draw=none, forget plot]
    table[x=iteration, y=ppo_lo] {figures/baselines_swimmer.dat};
\addplot[name path=ppoihi, draw=none, forget plot]
    table[x=iteration, y=ppo_hi] {figures/baselines_swimmer.dat};
\addplot[colorC, opacity=0.15, forget plot] fill between[of=ppoilo and ppoihi];
\addplot[colorC, line width=0.9pt]
    table[x=iteration, y=ppo_mean] {figures/baselines_swimmer.dat};
\addlegendentry{FedPPO};
\end{axis}
\end{tikzpicture}

    \caption{Swimmer-v4}
  \end{subfigure}
  \hfill
  \begin{subfigure}[b]{0.49\textwidth}
    %% Figure 3 analogue, Hopper-v4: three methods at N=8, mean over five seeds,
%% band spans min to max seed.
%% Data: figures/baselines_hopper.dat (scripts/export_figures.py).
\begin{tikzpicture}
\begin{axis}[
    width=0.96\linewidth,
    height=0.78\linewidth,
    xlabel={Iteration},
    xmin=1, xmax=2000,
    xtick distance=500,
    ylabel={Average batch return},
    ymin=0, ymax=3600,
    ytick distance=900,
    scaled y ticks=base 10:-3,
    ymajorgrids=true, xmajorgrids=true, grid style=dashed,
    tick label style={font=\footnotesize},
    label style={font=\footnotesize},
    legend style={font=\footnotesize, cells={anchor=west},
                  anchor=south east, at={(0.99,0.02)}},
]

\addplot[name path=admmilo, draw=none, forget plot]
    table[x=iteration, y=admm_lo] {figures/baselines_hopper.dat};
\addplot[name path=admmihi, draw=none, forget plot]
    table[x=iteration, y=admm_hi] {figures/baselines_hopper.dat};
\addplot[colorA, opacity=0.15, forget plot] fill between[of=admmilo and admmihi];
\addplot[colorA, line width=0.9pt, forget plot]
    table[x=iteration, y=admm_mean] {figures/baselines_hopper.dat};

\addplot[name path=stdilo, draw=none, forget plot]
    table[x=iteration, y=std_lo] {figures/baselines_hopper.dat};
\addplot[name path=stdihi, draw=none, forget plot]
    table[x=iteration, y=std_hi] {figures/baselines_hopper.dat};
\addplot[colorB, opacity=0.15, forget plot] fill between[of=stdilo and stdihi];
\addplot[colorB, line width=0.9pt, forget plot]
    table[x=iteration, y=std_mean] {figures/baselines_hopper.dat};

\addplot[name path=ppoilo, draw=none, forget plot]
    table[x=iteration, y=ppo_lo] {figures/baselines_hopper.dat};
\addplot[name path=ppoihi, draw=none, forget plot]
    table[x=iteration, y=ppo_hi] {figures/baselines_hopper.dat};
\addplot[colorC, opacity=0.15, forget plot] fill between[of=ppoilo and ppoihi];
\addplot[colorC, line width=0.9pt, forget plot]
    table[x=iteration, y=ppo_mean] {figures/baselines_hopper.dat};
\end{axis}
\end{tikzpicture}

    \caption{Hopper-v4}
  \end{subfigure}

  \vspace{0.6em}
  \begin{subfigure}[b]{0.49\textwidth}
    %% Figure 3 analogue, Humanoid-v4: three methods at N=8, mean over five seeds,
%% band spans min to max seed.
%% Data: figures/baselines_humanoid.dat (scripts/export_figures.py).
\begin{tikzpicture}
\begin{axis}[
    width=0.96\linewidth,
    height=0.78\linewidth,
    xlabel={Iteration},
    xmin=1, xmax=3000,
    xtick distance=1000,
    ylabel={Average batch return},
    ymin=0, ymax=8000,
    ytick distance=2000,
    scaled y ticks=base 10:-3,
    ymajorgrids=true, xmajorgrids=true, grid style=dashed,
    tick label style={font=\footnotesize},
    label style={font=\footnotesize},
    legend style={font=\footnotesize, cells={anchor=west},
                  anchor=south east, at={(0.99,0.02)}},
]

\addplot[name path=admmilo, draw=none, forget plot]
    table[x=iteration, y=admm_lo] {figures/baselines_humanoid.dat};
\addplot[name path=admmihi, draw=none, forget plot]
    table[x=iteration, y=admm_hi] {figures/baselines_humanoid.dat};
\addplot[colorA, opacity=0.15, forget plot] fill between[of=admmilo and admmihi];
\addplot[colorA, line width=0.9pt, forget plot]
    table[x=iteration, y=admm_mean] {figures/baselines_humanoid.dat};

\addplot[name path=stdilo, draw=none, forget plot]
    table[x=iteration, y=std_lo] {figures/baselines_humanoid.dat};
\addplot[name path=stdihi, draw=none, forget plot]
    table[x=iteration, y=std_hi] {figures/baselines_humanoid.dat};
\addplot[colorB, opacity=0.15, forget plot] fill between[of=stdilo and stdihi];
\addplot[colorB, line width=0.9pt, forget plot]
    table[x=iteration, y=std_mean] {figures/baselines_humanoid.dat};

\addplot[name path=ppoilo, draw=none, forget plot]
    table[x=iteration, y=ppo_lo] {figures/baselines_humanoid.dat};
\addplot[name path=ppoihi, draw=none, forget plot]
    table[x=iteration, y=ppo_hi] {figures/baselines_humanoid.dat};
\addplot[colorC, opacity=0.15, forget plot] fill between[of=ppoilo and ppoihi];
\addplot[colorC, line width=0.9pt, forget plot]
    table[x=iteration, y=ppo_mean] {figures/baselines_humanoid.dat};
\end{axis}
\end{tikzpicture}

    \caption{Humanoid-v4}
  \end{subfigure}
  \caption{The three methods at $N=8$ against iteration. Mean over five seeds,
  band from minimum to maximum seed, smoothed as in
  Figure~\ref{fig:convergence}. Colours as in the legend of panel (a).}
  \label{fig:baselines-iter}
\end{figure}

\begin{figure}[t]
  \centering
  \begin{subfigure}[b]{0.49\textwidth}
    %% Figure 3 analogue, Swimmer-v4: three methods at N=8, mean over five seeds,
%% band spans min to max seed.
%% x is cumulative uplink across the federation (64 bits per value).
%% Data: figures/baselines_swimmer.dat (scripts/export_figures.py).
\begin{tikzpicture}
\begin{axis}[
    width=0.96\linewidth,
    height=0.78\linewidth,
    xlabel={Cumulative uplink (bytes)},
    xmode=log,
    ylabel={Average batch return},
    ymin=0, ymax=140,
    ytick distance=35,
    ymajorgrids=true, xmajorgrids=true, grid style=dashed,
    tick label style={font=\footnotesize},
    label style={font=\footnotesize},
    legend style={font=\footnotesize, cells={anchor=west},
                  anchor=north west, at={(0.02,0.98)}},
]

\addplot[name path=admmblo, draw=none, forget plot]
    table[x=admm_bytes, y=admm_lo] {figures/baselines_swimmer.dat};
\addplot[name path=admmbhi, draw=none, forget plot]
    table[x=admm_bytes, y=admm_hi] {figures/baselines_swimmer.dat};
\addplot[colorA, opacity=0.15, forget plot] fill between[of=admmblo and admmbhi];
\addplot[colorA, line width=0.9pt]
    table[x=admm_bytes, y=admm_mean] {figures/baselines_swimmer.dat};
\addlegendentry{FedNPG-ADMM};

\addplot[name path=stdblo, draw=none, forget plot]
    table[x=std_bytes, y=std_lo] {figures/baselines_swimmer.dat};
\addplot[name path=stdbhi, draw=none, forget plot]
    table[x=std_bytes, y=std_hi] {figures/baselines_swimmer.dat};
\addplot[colorB, opacity=0.15, forget plot] fill between[of=stdblo and stdbhi];
\addplot[colorB, line width=0.9pt]
    table[x=std_bytes, y=std_mean] {figures/baselines_swimmer.dat};
\addlegendentry{Standard FedNPG};

\addplot[name path=ppoblo, draw=none, forget plot]
    table[x=ppo_bytes, y=ppo_lo] {figures/baselines_swimmer.dat};
\addplot[name path=ppobhi, draw=none, forget plot]
    table[x=ppo_bytes, y=ppo_hi] {figures/baselines_swimmer.dat};
\addplot[colorC, opacity=0.15, forget plot] fill between[of=ppoblo and ppobhi];
\addplot[colorC, line width=0.9pt]
    table[x=ppo_bytes, y=ppo_mean] {figures/baselines_swimmer.dat};
\addlegendentry{FedPPO};
\end{axis}
\end{tikzpicture}

    \caption{Swimmer-v4}
  \end{subfigure}
  \hfill
  \begin{subfigure}[b]{0.49\textwidth}
    %% Figure 3 analogue, Hopper-v4: three methods at N=8, mean over five seeds,
%% band spans min to max seed.
%% x is cumulative uplink across the federation (64 bits per value).
%% Data: figures/baselines_hopper.dat (scripts/export_figures.py).
\begin{tikzpicture}
\begin{axis}[
    width=0.96\linewidth,
    height=0.78\linewidth,
    xlabel={Cumulative uplink (bytes)},
    xmode=log,
    ylabel={Average batch return},
    ymin=0, ymax=3600,
    ytick distance=900,
    scaled y ticks=base 10:-3,
    ymajorgrids=true, xmajorgrids=true, grid style=dashed,
    tick label style={font=\footnotesize},
    label style={font=\footnotesize},
    legend style={font=\footnotesize, cells={anchor=west},
                  anchor=south east, at={(0.99,0.02)}},
]

\addplot[name path=admmblo, draw=none, forget plot]
    table[x=admm_bytes, y=admm_lo] {figures/baselines_hopper.dat};
\addplot[name path=admmbhi, draw=none, forget plot]
    table[x=admm_bytes, y=admm_hi] {figures/baselines_hopper.dat};
\addplot[colorA, opacity=0.15, forget plot] fill between[of=admmblo and admmbhi];
\addplot[colorA, line width=0.9pt, forget plot]
    table[x=admm_bytes, y=admm_mean] {figures/baselines_hopper.dat};

\addplot[name path=stdblo, draw=none, forget plot]
    table[x=std_bytes, y=std_lo] {figures/baselines_hopper.dat};
\addplot[name path=stdbhi, draw=none, forget plot]
    table[x=std_bytes, y=std_hi] {figures/baselines_hopper.dat};
\addplot[colorB, opacity=0.15, forget plot] fill between[of=stdblo and stdbhi];
\addplot[colorB, line width=0.9pt, forget plot]
    table[x=std_bytes, y=std_mean] {figures/baselines_hopper.dat};

\addplot[name path=ppoblo, draw=none, forget plot]
    table[x=ppo_bytes, y=ppo_lo] {figures/baselines_hopper.dat};
\addplot[name path=ppobhi, draw=none, forget plot]
    table[x=ppo_bytes, y=ppo_hi] {figures/baselines_hopper.dat};
\addplot[colorC, opacity=0.15, forget plot] fill between[of=ppoblo and ppobhi];
\addplot[colorC, line width=0.9pt, forget plot]
    table[x=ppo_bytes, y=ppo_mean] {figures/baselines_hopper.dat};
\end{axis}
\end{tikzpicture}

    \caption{Hopper-v4}
  \end{subfigure}

  \vspace{0.6em}
  \begin{subfigure}[b]{0.49\textwidth}
    %% Figure 3 analogue, Humanoid-v4: three methods at N=8, mean over five seeds,
%% band spans min to max seed.
%% x is cumulative uplink across the federation (64 bits per value).
%% Data: figures/baselines_humanoid.dat (scripts/export_figures.py).
\begin{tikzpicture}
\begin{axis}[
    width=0.96\linewidth,
    height=0.78\linewidth,
    xlabel={Cumulative uplink (bytes)},
    xmode=log,
    ylabel={Average batch return},
    ymin=0, ymax=8000,
    ytick distance=2000,
    scaled y ticks=base 10:-3,
    ymajorgrids=true, xmajorgrids=true, grid style=dashed,
    tick label style={font=\footnotesize},
    label style={font=\footnotesize},
    legend style={font=\footnotesize, cells={anchor=west},
                  anchor=south east, at={(0.99,0.02)}},
]

\addplot[name path=admmblo, draw=none, forget plot]
    table[x=admm_bytes, y=admm_lo] {figures/baselines_humanoid.dat};
\addplot[name path=admmbhi, draw=none, forget plot]
    table[x=admm_bytes, y=admm_hi] {figures/baselines_humanoid.dat};
\addplot[colorA, opacity=0.15, forget plot] fill between[of=admmblo and admmbhi];
\addplot[colorA, line width=0.9pt, forget plot]
    table[x=admm_bytes, y=admm_mean] {figures/baselines_humanoid.dat};

\addplot[name path=stdblo, draw=none, forget plot]
    table[x=std_bytes, y=std_lo] {figures/baselines_humanoid.dat};
\addplot[name path=stdbhi, draw=none, forget plot]
    table[x=std_bytes, y=std_hi] {figures/baselines_humanoid.dat};
\addplot[colorB, opacity=0.15, forget plot] fill between[of=stdblo and stdbhi];
\addplot[colorB, line width=0.9pt, forget plot]
    table[x=std_bytes, y=std_mean] {figures/baselines_humanoid.dat};

\addplot[name path=ppoblo, draw=none, forget plot]
    table[x=ppo_bytes, y=ppo_lo] {figures/baselines_humanoid.dat};
\addplot[name path=ppobhi, draw=none, forget plot]
    table[x=ppo_bytes, y=ppo_hi] {figures/baselines_humanoid.dat};
\addplot[colorC, opacity=0.15, forget plot] fill between[of=ppoblo and ppobhi];
\addplot[colorC, line width=0.9pt, forget plot]
    table[x=ppo_bytes, y=ppo_mean] {figures/baselines_humanoid.dat};
\end{axis}
\end{tikzpicture}

    \caption{Humanoid-v4}
  \end{subfigure}
  \caption{The same runs against cumulative uplink traffic across the
  federation, on a logarithmic axis, at 64 bits per transmitted value.
  Standard-average FedNPG sits three to six orders of magnitude to the right
  because each agent uplinks its $d \times d$ Fisher matrix every round.
  Colours as in the legend of panel (a).}
  \label{fig:baselines-bytes}
\end{figure}

Table~\ref{tab:baselines} and Figure~\ref{fig:baselines-iter} compare the
methods per iteration. On Swimmer the three are indistinguishable: they finish
within three points of each other, and the environment saturates near $128$
for every method, as it does for every agent count in
Section~\ref{sec:reproduction}. On Hopper, FedNPG-ADMM and standard-average
FedNPG agree to within the seed spread ($2656 \pm 286$ against
$2494 \pm 677$), which is the paper's central empirical claim: the ADMM
approximation costs no reward relative to the exact aggregation it replaces,
here at $\num{9219}$ times less uplink per round. FedPPO finishes well below
both, at $1571 \pm 588$.

Humanoid departs from that pattern. Standard-average FedNPG, which computes
exactly the pooled natural-gradient direction that FedNPG-ADMM approximates
with a single ADMM round, finishes \emph{below} it, $4699 \pm 298$ against
$6329 \pm 334$, a gap several times the seed spread. We have not established
why. One candidate is that the single-round iterate is not a pure
approximation: through the dual variables and the $\rho\, \by^{k-1}$ term of
line~8, $\by^k$ carries information from previous iterations, which acts
somewhat like momentum.

Figure~\ref{fig:baselines-bytes} plots the same curves against uplink traffic. Against standard
FedNPG the comparison is not close on any environment. Against FedPPO, which
transmits half as much per round, the outcome depends on the environment. At
the total traffic FedPPO uses over its whole run, FedNPG-ADMM has completed
half as many iterations: on Swimmer both stand at about $128$; on Hopper
FedNPG-ADMM is ahead ($2471$ against $1763$ on the smoothed curves); on
Humanoid FedPPO is ahead ($5973$ against $5130$). Per byte, then, FedNPG-ADMM
is not uniformly the better method. Its advantage over FedPPO is in reward per
round, which it holds on Hopper and Humanoid, and the reduction it delivers is
the one the paper claims, against the $\Oh(d^2)$ method it was designed to
replace.
